% ============================================================
%  Capítulo: Interfaces de línea de comandos, entornos virtuales
%            y scripting de intérprete para Windows y Linux
% ============================================================
%  Dependencias (agregar al preámbulo si aún no están presentes):
%    \usepackage{listings}
%    \usepackage{xcolor}
%    \usepackage{booktabs}
%    \usepackage{hyperref}
%    \usepackage{tcolorbox}
%    \usepackage{fontawesome5}   % opcional, para iconos
% ============================================================


% ============================================================
\chapter{Interfaces de línea de comandos, entornos virtuales
         y scripting de intérprete}
\label{ch:cli-scripting}
% ============================================================

% ============================================================
\section{¿Qué es un archivo?}
\label{sec:what-is-a-file}
% ============================================================

\subsection{Archivos y texto}

Un archivo es una secuencia nombrada de bytes almacenada en un dispositivo de
almacenamiento. Desde la perspectiva de la computadora, todos los archivos son
secuencias de bytes. Lo que distingue a un script de Python de una fotografía
no es el mecanismo de almacenamiento sino la interpretación: un editor de
texto interpreta los bytes como caracteres; un visor de imágenes los
interpreta como valores de píxeles. Un archivo de Python es texto plano: cada
carácter que usted escribe se almacena como un byte, y el intérprete de
Python lee esos bytes como instrucciones. Por esta razón un documento de Word
no puede servir como un programa de Python; Word almacena bytes de formato
invisibles junto con el texto.

\subsection{El sistema de archivos: carpetas y rutas}

Cada archivo reside dentro de una carpeta (también llamada directorio). Las
carpetas pueden contener otras carpetas, formando un árbol. La ubicación de
cualquier archivo puede describirse mediante una \emph{ruta}: una secuencia
de nombres de carpetas separados por un delimitador, que termina con el
nombre del archivo.

\begin{table}[htbp]
\centering
\caption{Sintaxis de rutas en cada sistema operativo.}
\label{tab:path-syntax}
\begin{tabular}{p{3.5cm} p{9cm}}
\toprule
\textbf{Sistema operativo} & \textbf{Ruta de ejemplo} \\
\midrule
Windows & \texttt{C:\textbackslash Users\textbackslash Alice\textbackslash Documents\textbackslash hello.py} \\
macOS   & \texttt{/Users/alice/Documents/hello.py} \\
Linux   & \texttt{/home/alice/Documents/hello.py} \\
\bottomrule
\end{tabular}
\end{table}

Windows utiliza barras invertidas y una letra de unidad; macOS y Linux
utilizan barras diagonales y parten de una raíz \texttt{/}.

\subsection{El directorio de trabajo}

En cualquier momento, la línea de comandos está «dentro» de una carpeta en
particular, llamada el \emph{directorio de trabajo actual} (CWD, por sus
siglas en inglés). Los comandos que hacen referencia a archivos sin una ruta
completa se interpretan relativos a esa carpeta. Cuando usted escribe
\texttt{python hello.py}, el intérprete busca \texttt{hello.py} en el CWD. Si
el archivo se encuentra en otra ubicación, el comando falla. Comprender el
CWD es la fuente más común de confusión para las personas que recién empiezan
a usar la línea de comandos y la causa raíz de la mayoría de los errores de
«archivo no encontrado».

La línea de comandos proporciona los comandos para navegar el sistema de
archivos, inspeccionar su contenido y ejecutar archivos; estos se presentan
en la siguiente sección.

% ============================================================

\begin{quote}
\textit{«La línea de comandos es la interfaz más poderosa jamás inventada.
También es la más intimidante. Este capítulo pretende cambiar eso.»}
\end{quote}

% ============================================================
\section{Introducción: ¿qué es una interfaz de línea de comandos?}
\label{sec:cli-intro}
% ============================================================

Cuando usted utiliza una computadora, seguramente interactúa con ella a
través de una \textbf{interfaz gráfica de usuario} (GUI, por sus siglas en
inglés): ventanas, iconos, menús y botones sobre los que hace clic con el
ratón. Las GUI son intuitivas y visualmente ricas, pero tienen una
limitación fundamental: cada acción posible debe haber sido anticipada e
implementada como un botón o elemento de menú por la persona que desarrolló
el software. Si usted necesita hacer algo para lo que la GUI no fue diseñada,
queda bloqueado.

Una \textbf{interfaz de línea de comandos} (CLI) resuelve esto al permitirle
comunicarse con la computadora mediante \emph{texto}. Usted escribe un
comando; la computadora lo ejecuta y responde con texto. Esto no es un paso
atrás; es un modo de interacción diferente y, en muchos contextos, más
poderoso.

\subsection{Una analogía concreta}

Piense en la diferencia entre un restaurante con un menú fijo y un
restaurante donde usted puede describir exactamente lo que desea al chef.
Una GUI es el menú fijo: rápida, amigable y suficientemente buena para la
mayoría de los clientes. Una CLI es hablar directamente con el chef: requiere
más conocimiento de los ingredientes, pero los resultados pueden ser
precisamente lo que usted necesita, y puede combinarlos de formas que el
menú nunca anticipó.

\subsection{Por qué las personas que desarrollan y las que investigan usan la CLI}

\begin{itemize}
  \item \textbf{Automatización.} Una secuencia de comandos puede guardarse
        en un archivo (llamado \emph{script}) y reproducirse a la perfección
        cada vez.
  \item \textbf{Precisión.} Cada opción y parámetro es explícito y legible.
  \item \textbf{Acceso remoto.} Los servidores en la nube rara vez disponen
        de un escritorio gráfico; la CLI es a menudo la única vía de entrada.
  \item \textbf{Velocidad.} Cambiar el nombre de 10{,}000 archivos requiere
        un solo comando; hacerlo a mano en una GUI podría tomar horas.
  \item \textbf{Reproducibilidad.} Usted puede compartir un script con una
        persona colega y saber que obtendrá exactamente el mismo resultado.
\end{itemize}

\subsection{Los dos mundos principales de la CLI}

Este capítulo se centra en dos entornos de CLI:

\begin{description}
  \item[Command Prompt / PowerShell de Windows] Las herramientas CLI nativas
        en Microsoft Windows. Los archivos por lotes (\texttt{.bat}) y los
        scripts de PowerShell (\texttt{.ps1}) automatizan tareas aquí.
  \item[Bash (Linux / macOS)] El intérprete de comandos dominante en
        sistemas Linux y macOS. Los scripts de intérprete (\texttt{.sh})
        automatizan tareas aquí.
\end{description}

Estos dos mundos utilizan sintaxis diferentes, convenciones de ruta diferentes
(barra invertida \texttt{\textbackslash} vs.\ barra diagonal \texttt{/}) y
filosofías de scripting diferentes. Un objetivo clave de este capítulo es
ayudarle a desplazarse con fluidez entre ellos.

\subsection{Sus primeros comandos}

La mejor manera de desmitificar la CLI es probarla. La
Tabla~\ref{tab:first-commands} muestra comandos equivalentes en ambos
entornos.

\begin{table}[h]
\centering
\caption{Comandos CLI equivalentes: Windows vs.\ Linux/macOS}
\label{tab:first-commands}
\begin{tabular}{@{}lll@{}}
\toprule
\textbf{Tarea} & \textbf{Windows (CMD)} & \textbf{Linux / macOS (Bash)} \\
\midrule
Mostrar directorio actual   & \texttt{cd}              & \texttt{pwd}             \\
Listar archivos             & \texttt{dir}             & \texttt{ls -la}          \\
Cambiar de directorio       & \texttt{cd foldername}   & \texttt{cd foldername}   \\
Subir un nivel              & \texttt{cd ..}           & \texttt{cd ..}           \\
Crear una carpeta           & \texttt{mkdir myfolder}  & \texttt{mkdir myfolder}  \\
Eliminar un archivo         & \texttt{del file.txt}    & \texttt{rm file.txt}     \\
Copiar un archivo           & \texttt{copy a.txt b.txt}& \texttt{cp a.txt b.txt}  \\
Imprimir texto en pantalla  & \texttt{echo Hello}      & \texttt{echo Hello}      \\
Limpiar la pantalla         & \texttt{cls}             & \texttt{clear}           \\
\bottomrule
\end{tabular}
\end{table}

\begin{tipbox}
En ambos entornos, presionar la tecla de \textbf{flecha arriba} recupera sus
comandos anteriores, y presionar \textbf{Tab} autocompleta nombres de
archivos y directorios. Solamente estos dos hábitos acelerarán
notablemente su trabajo en la línea de comandos.
\end{tipbox}

% ============================================================
\section{Dos mundos: Windows y Linux/Mac}
\label{sec:two-worlds}
% ============================================================

Antes de escribir cualquier script, conviene comprender las diferencias
filosóficas entre la forma en que Windows y Linux manejan la línea de
comandos. Estas diferencias no son arbitrarias; reflejan décadas de
historias de diseño divergentes.

\subsection{Rutas y separadores}

Windows utiliza \emph{barras invertidas} y letras de unidad para describir
ubicaciones de archivos:
\begin{lstlisting}[style=windowsStyle,escapeinside={(*@}{@*)}]
C:\Users\DrG\Documents\(*@\ProjectName@*)\main.py
\end{lstlisting}

Linux utiliza \emph{barras diagonales} y un único árbol de archivos
unificado que tiene como raíz \texttt{/}:
\begin{lstlisting}[style=linuxStyle,escapeinside={(*@}{@*)}]
/home/drg/(*@\ProjectName@*)/main.py
\end{lstlisting}

\subsection{Scripts y ejecutabilidad}

En Windows, la extensión del archivo determina si un archivo es ejecutable.
Los archivos que terminan en \texttt{.bat}, \texttt{.cmd} o \texttt{.exe}
se tratan como programas.

En Linux, la extensión es irrelevante. La ejecutabilidad es un
\emph{permiso} almacenado en el archivo mismo. Debe otorgarse
explícitamente permiso de ejecución a un archivo antes de que el sistema
lo ejecute. Esto se hace con el comando \texttt{chmod}, que cubrimos en
la Sección~\ref{sec:chmod}.

\subsection{Variables de entorno y PATH}

Ambos sistemas utilizan una variable de entorno llamada \texttt{PATH} para
encontrar programas. Cuando usted escribe un comando, el intérprete de
comandos busca en cada directorio listado en \texttt{PATH} en orden y
ejecuta la primera coincidencia que encuentra.

\begin{itemize}
  \item \textbf{Windows:} PATH se administra a través de Propiedades del
        sistema $\to$ Variables de entorno, o mediante el comando
        \texttt{setx}.
  \item \textbf{Linux:} PATH se configura en archivos de configuración como
        \texttt{\textasciitilde/.bashrc} (para el intérprete Bash) o
        \texttt{\textasciitilde/.zshrc} (para Zsh), los cuales se leen cada
        vez que se inicia una nueva sesión de terminal.
\end{itemize}

% ============================================================
\section{WSL: un entorno Linux dentro de Windows}
\label{sec:wsl}
% ============================================================

Para las personas que utilizan Windows y desean el poder del scripting de
Linux sin abandonar su entorno familiar, Microsoft proporciona el
\textbf{Subsistema de Windows para Linux} (WSL). WSL no es un emulador ni
una máquina virtual en el sentido tradicional; es una capa de compatibilidad
integrada en Windows que ejecuta un núcleo Linux real junto al núcleo de
Windows.

\subsection{Qué le ofrece WSL}

\begin{itemize}
  \item Un intérprete Bash completo (o Zsh, Fish, etc.)
  \item Un sistema de archivos completo de Ubuntu (u otra distribución)
  \item Acceso a miles de herramientas de línea de comandos de Linux a
        través de \texttt{apt}
  \item La capacidad de leer y escribir archivos de Windows desde Linux
        (su unidad \texttt{C:} de Windows aparece en \texttt{/mnt/c/})
  \item Acceso a GPU para cargas de trabajo de aprendizaje automático
\end{itemize}

\subsection{Instalación de WSL}

Abra PowerShell o CMD \emph{como Administrador} y ejecute:

\begin{lstlisting}[style=windowsStyle, caption={Instalación de WSL desde PowerShell}]
wsl --install
\end{lstlisting}

Esto instala WSL 2 y la distribución Ubuntu de forma predeterminada.
Después de reiniciar, se abrirá una terminal de Ubuntu y le pedirá que cree
un nombre de usuario y una contraseña. Esa cuenta es independiente de su
inicio de sesión de Windows.

\begin{notebox}
A lo largo del resto de este capítulo, todas las instrucciones para Linux
se aplican por igual a WSL en Windows y a una máquina Linux nativa.
Trátelas como el mismo entorno.
\end{notebox}

\subsection{La solución elegante: Visual Studio Code y la extensión WSL}
\label{subsec:vscode-wsl}

Uno de los aspectos más poderosos de WSL es su integración con
\textbf{Visual Studio Code} (VS Code), el editor de código gratuito de
Microsoft. La instalación de la \textbf{extensión WSL} en VS Code crea un
puente sin interrupciones entre su escritorio de Windows y su entorno de
Linux:

\begin{itemize}
  \item Usted abre VS Code desde la terminal de WSL escribiendo
        \texttt{code .}
  \item VS Code ejecuta su interfaz en Windows pero ejecuta todo el código,
        los comandos de terminal, las extensiones y los depuradores
        \emph{dentro de Linux}
  \item El explorador de archivos muestra su sistema de archivos de Linux
  \item La terminal integrada es un intérprete Bash
  \item Los entornos de Python, los linters y los formateadores se resuelven
        a sus versiones de Linux
\end{itemize}

Esto significa que usted nunca tiene que elegir entre el ecosistema de
scripting de Linux y la experiencia de escritorio de Windows: obtiene ambos
simultáneamente. A partir de este punto, cuando digamos «abra VS Code»,
nos referiremos a iniciarlo desde la terminal de WSL con \texttt{code .}

\begin{tipbox}
Para instalar la extensión WSL, abra VS Code, presione
\texttt{Ctrl+Shift+X} para abrir el panel de Extensiones y busque
«WSL» (editor: Microsoft). Un clic la instala.
\end{tipbox}

% ============================================================
% Fin del capítulo
% ============================================================
